\subsection{Diagonalization, Commutants, and Simultaneous Forms}

\begin{definition} \label{def:diagonalizable-operator} \hlt{Diagonalizable Operator}\\
Let $\mathcal{V}$ be a finite-dimensional vector space over $\mathbb{F}$. An operator $T\in\LL(\mathcal{V})$ is \hlt{diagonalizable} iff $(\exists\text{ a basis }\beta\text{ of }\mathcal{V})\,[T]_{\beta\leftarrow\beta}$ is diagonal.
\end{definition}

\begin{proposition} \label{prop:eigenspaces-for-distinct-eigenvalues-form-direct-sum} \hlt{Eigenspaces for Distinct Eigenvalues Form a Direct Sum}\\
Let $\mathcal{V}$ be finite-dimensional, $T\in\LL(\mathcal{V})$, and let $\lambda_1,\ldots,\lambda_m$ be distinct eigenvalues of $T$, where $m\geq1$. The following hold:
\begin{enumerate}[label=(\roman*)]
\setlength{\itemsep}{0pt}
\item The sum $E_{\lambda_1}(T)+\cdots+E_{\lambda_m}(T)$ is direct.
\item $\sum\limits_{j=1}^{m}\dim E_{\lambda_j}(T)\leq\dim(\mathcal{V})$.
\end{enumerate}
\begin{proof}
\begin{enumerate}[label=(\roman*)]
\setlength{\itemsep}{0pt}
\item Suppose $(\mathbf{u}_j)_{j=1}^{m}\in\prod\limits_{j=1}^{m}E_{\lambda_j}(T)$ and $\sum\limits_{j=1}^{m}\mathbf{u}_j=\mathbf{0}$. Set $J\defeq\{j\in\{1,\ldots,m\}:\mathbf{u}_j\neq\mathbf{0}\}$.\\
If $J\neq\emptyset$, then $(\mathbf{u}_j)_{j\in J}$ is linearly independent by Theorem \ref{thm:distinct-eigenvalues-independent-eigenvectors}. This contradicts its zero relation with coefficients $1$.\\
Thus $J=\emptyset$, so the sum is direct by Proposition \ref{prop:internal-direct-sum-family-criterion}(ii),(i).
\item $\dim(E_{\lambda_1}(T)\oplus\cdots\oplus E_{\lambda_m}(T))=\sum\limits_{j=1}^{m}\dim E_{\lambda_j}(T)$ by Theorem \ref{thm:direct-sum-iff-dimensions-add}(i). The left side is at most $\dim(\mathcal{V})$ by Proposition \ref{prop:dimension-bounds-and-basis-criteria}(i).
\end{enumerate}
\hlt{Structural Commentary}: Removing zero components reduces every mixed relation to distinct-eigenvalue independence.\\
\hlt{Key Insight}: Distinct-eigenvalue independence forbids cross-eigenspace cancellation.
\end{proof}
\end{proposition}

\begin{theorem} \label{thm:conditions-equivalent-to-diagonalizability} \hlt{Conditions Equivalent to Diagonalizability}\\
Let $\mathcal{V}$ be finite-dimensional, $T\in\LL(\mathcal{V})$, and $n\defeq\dim(\mathcal{V})$.\\
Let $(\lambda_1,\ldots,\lambda_m)$ list the distinct elements of $\operatorname{spec}(T)$, where $m\geq0$. The following are equivalent:
\begin{enumerate}[label=(\roman*)]
\setlength{\itemsep}{0pt}
\item $T$ is diagonalizable;
\item $\mathcal{V}$ has a basis consisting of eigenvectors of $T$;
\item $(\exists\,(\mathcal{U}_j)_{j=1}^{n}\in\operatorname{Inv}(T)^n)\,[(\forall j\in\{1,\ldots,n\})\,\dim(\mathcal{U}_j)=1\land\mathcal{V}=\bigoplus\limits_{j=1}^{n}\mathcal{U}_j]$;
\item $\mathcal{V}=\bigoplus\limits_{j=1}^{m}E_{\lambda_j}(T)$;
\item $\dim(\mathcal{V})=\sum\limits_{j=1}^{m}\dim E_{\lambda_j}(T)$.
\end{enumerate}
\begin{proof}
\begin{enumerate}[label=(\roman*)]
\setlength{\itemsep}{0pt}
\item[$(i)$] $\Rightarrow (ii)$ Choose a basis $\beta=(\mathbf{v}_1,\ldots,\mathbf{v}_n)$ such that $[T]_{\beta\leftarrow\beta}$ is diagonal.\\
Writing its $j$th column as $a_j\mathbf{e}_j$, $T\mathbf{v}_j=a_j\mathbf{v}_j$ by the matrix representation. Thus $\beta$ is an eigenbasis.
\item[$(ii)$] $\Rightarrow (iii)$ Let $(\mathbf{v}_1,\ldots,\mathbf{v}_n)$ be an eigenbasis and set $\mathcal{U}_j\defeq\operatorname{span}(\mathbf{v}_j)$.\\
Each $\mathcal{U}_j$ is one-dimensional and $T$-invariant by Proposition \ref{prop:one-dimensional-invariant-subspaces-and-eigenvectors}. The equality $\mathcal{V}=\sum\limits_{j=1}^{n}\mathcal{U}_j$ holds by spanning.\\
Every zero relation across the summands is a basis relation, so the sum is direct by Proposition \ref{prop:internal-direct-sum-family-criterion}(ii),(i).
\item[$(iii)$] $\Rightarrow (iv)$ Choose $(\mathbf{v}_j)_{j=1}^{n}\in\prod\limits_{j=1}^{n}(\mathcal{U}_j\setminus\{\mathbf{0}\})$. Each $\mathbf{v}_j$ is an eigenvector by Proposition \ref{prop:one-dimensional-invariant-subspaces-and-eigenvectors}.\\
Thus every $\mathcal{U}_j$ lies in an eigenspace of $T$, and the listed eigenspaces span $\mathcal{V}$.\\
If $m=0$ and $n\geq1$, then $\mathbf{v}_1$ contradicts $\operatorname{spec}(T)=\emptyset$; hence $m=0\Rightarrow n=0$.\\
If $m=0$, the asserted empty direct sum is $\mathcal{V}=\{\mathbf{0}\}$. If $m\geq1$, the sum is direct by Proposition \ref{prop:eigenspaces-for-distinct-eigenvalues-form-direct-sum}(i).
\item[$(iv)$] $\Rightarrow (v)$ If $m=0$, condition (iv) states $\mathcal{V}=\{\mathbf{0}\}$, so both sides of (v) vanish. If $m\geq1$, the dimensions add by Theorem \ref{thm:direct-sum-iff-dimensions-add}(i).
\item[$(v)$] $\Rightarrow (i)$ If $m=0$, then $n=0$ by condition (v), and the empty basis diagonalizes $T$. Suppose $m\geq1$.\\
Choose bases $(\beta_j)_{j=1}^{m}$ with $\beta_j$ a basis of $E_{\lambda_j}(T)$, and let $\beta$ be their concatenation.\\
The eigenspace sum is direct by Proposition \ref{prop:eigenspaces-for-distinct-eigenvalues-form-direct-sum}(i). Hence $\beta$ is a basis of that sum by Theorem \ref{thm:direct-sum-iff-dimensions-add}(ii).\\
Its length is $n$ by (v), so it is a basis of $\mathcal{V}$ by Proposition \ref{prop:full-dimension-basis-criteria}(ii).\\
Each member of $\beta$ is an eigenvector; hence $[T]_{\beta\leftarrow\beta}$ is diagonal, and $T$ is diagonalizable.
\end{enumerate}
\hlt{Structural Commentary}: The implication cycle converts diagonal coordinates into invariant lines and exhaustive eigenspace dimension.\\
\hlt{Key Insight}: Diagonalizability is exactly exhaustion by distinct eigenspaces.
\end{proof}
\end{theorem}

\begin{corollary} \label{cor:enough-eigenvalues-implies-diagonalizable} \hlt{Enough Eigenvalues Imply Diagonalizability}\\
Let $\mathcal{V}$ be finite-dimensional and let $T\in\LL(\mathcal{V})$. If $\lvert\operatorname{spec}(T)\rvert=\dim(\mathcal{V})$, then $T$ is diagonalizable.
\begin{proof}
Set $n\defeq\dim(\mathcal{V})$. If $n=0$, the empty basis diagonalizes $T$. Suppose $n\geq1$ and list $\operatorname{spec}(T)$ as $\lambda_1,\ldots,\lambda_n$.\\
Choose $(\mathbf{v}_j)_{j=1}^{n}\in\prod\limits_{j=1}^{n}(E_{\lambda_j}(T)\setminus\{\mathbf{0}\})$. The list $(\mathbf{v}_1,\ldots,\mathbf{v}_n)$ is linearly independent by Theorem \ref{thm:distinct-eigenvalues-independent-eigenvectors}.\\
Its length is $n$, so it is a basis by Proposition \ref{prop:full-dimension-basis-criteria}(ii). Thus $T$ is diagonalizable by Theorem \ref{thm:conditions-equivalent-to-diagonalizability}(ii),(i).\\
\hlt{Structural Commentary}: Distinct eigenvalues supply an independent eigenvector list whose cardinality reaches the ambient dimension.\\
\hlt{Key Insight}: A full-length independent eigenvector list is already a diagonalizing basis.
\end{proof}
\end{corollary}

\begin{perspective} \label{per:diagonalizability-viewpoint} \hlt{Diagonalization as Exhaustion by Eigenspaces}\\
Separation holds by Proposition \ref{prop:eigenspaces-for-distinct-eigenvalues-form-direct-sum}(i). Spanning is the remaining condition by Theorem \ref{thm:conditions-equivalent-to-diagonalizability}.\\
Diagonalization therefore records exhaustion of $\mathcal{V}$ by its distinct eigenspaces.
\end{perspective}

\begin{theorem} \label{thm:algebraic-spectral-resolution-diagonalizable} \hlt{Algebraic Spectral Resolution}\\
Let $\mathcal{V}$ be nonzero and finite-dimensional, $T\in\LL(\mathcal{V})$, and let $\lambda_1,\ldots,\lambda_m\in\mathbb{F}$ be distinct, where $m\geq1$.\\
Set $\Lambda\defeq\{\lambda_1,\ldots,\lambda_m\}$, and interpret an empty product below as $I_{\mathcal{V}}$. The following hold:
\begin{enumerate}[label=(\roman*)]
\setlength{\itemsep}{0pt}
\item $T$ is diagonalizable with $\operatorname{spec}(T)=\Lambda$ iff $(\exists\,(\rho_i)_{i=1}^{m}\in(\LL(\mathcal{V})\setminus\{0\})^m)$ such that $(\forall i\in\{1,\ldots,m\})\,\rho_i^2=\rho_i$,\\
$(\forall i,j\in\{1,\ldots,m\})\,[i\neq j\Rightarrow\rho_i\rho_j=0]$, $\sum\limits_{i=1}^{m}\rho_i=I_{\mathcal{V}}$, and $T=\sum\limits_{i=1}^{m}\lambda_i\rho_i$.
\item Every family $(\rho_i)_{i=1}^{m}$ satisfying the right side of part (i) obeys $(\forall i\in\{1,\ldots,m\})\,\operatorname{range}(\rho_i)=E_{\lambda_i}(T)$.
\item Every such family obeys $(\forall i\in\{1,\ldots,m\})\,\ker(\rho_i)=\bigoplus\limits_{j\in\{1,\ldots,m\}\setminus\{i\}}E_{\lambda_j}(T)$.
\item Every such family obeys $(\forall i\in\{1,\ldots,m\})\,\rho_i=\prod\limits_{j\in\{1,\ldots,m\}\setminus\{i\}}\frac{T-\lambda_jI_{\mathcal{V}}}{\lambda_i-\lambda_j}$.
\end{enumerate}
\begin{proof}
\begin{enumerate}[label=(\roman*)]
\setlength{\itemsep}{0pt}
\item[$(i)$] ($\Rightarrow$) The decomposition $\mathcal{V}=\bigoplus\limits_{i=1}^{m}E_{\lambda_i}(T)$ holds by Theorem \ref{thm:conditions-equivalent-to-diagonalizability}(i),(iv).\\
For $i\in\{1,\ldots,m\}$, set $\mathcal{W}_i\defeq\bigoplus\limits_{j\in\{1,\ldots,m\}\setminus\{i\}}E_{\lambda_j}(T)$ and $\rho_i\defeq\rho_{E_{\lambda_i}(T),\mathcal{W}_i}$.\\
Each $\rho_i$ is idempotent by Proposition \ref{prop:algebraic-projection-range-kernel-idempotent}(ii), and $\rho_i\neq0$ because $E_{\lambda_i}(T)\neq\{\mathbf{0}\}$.\\
If $\mathbf{x}=\sum\limits_{k=1}^{m}\mathbf{x}_k$ with $\mathbf{x}_k\in E_{\lambda_k}(T)$, then $(\forall i\in\{1,\ldots,m\})\,\rho_i\mathbf{x}=\mathbf{x}_i$ by the projection rule.\\
Thus $(\forall i,j\in\{1,\ldots,m\})\,[i\neq j\Rightarrow\rho_i\rho_j=0]$ and $\sum\limits_{i=1}^{m}\rho_i=I_{\mathcal{V}}$.\\
The identity $(\forall\mathbf{x}\in\mathcal{V})\,T\mathbf{x}=\sum\limits_{i=1}^{m}\lambda_i\mathbf{x}_i=\sum\limits_{i=1}^{m}\lambda_i\rho_i\mathbf{x}$ holds by the eigenspace identities and projection rule, so $T=\sum\limits_{i=1}^{m}\lambda_i\rho_i$.
\item[$(i)$] ($\Leftarrow$) Fix a family satisfying the right side, and set $\mathcal{V}_i\defeq\operatorname{range}(\rho_i)$ for $i\in\{1,\ldots,m\}$. The equality $\mathcal{V}=\sum\limits_{i=1}^{m}\mathcal{V}_i$ holds by $\sum\limits_{i=1}^{m}\rho_i=I_{\mathcal{V}}$.\\
Suppose $(\mathbf{x}_i)_{i=1}^{m}\in\prod\limits_{i=1}^{m}\mathcal{V}_i$ and $\sum\limits_{i=1}^{m}\mathbf{x}_i=\mathbf{0}$. Then $(\forall k\in\{1,\ldots,m\})\,\mathbf{x}_k=\rho_k\sum\limits_{i=1}^{m}\mathbf{x}_i=\mathbf{0}$ by idempotence and pairwise annihilation.\\
Hence $\mathcal{V}=\bigoplus\limits_{i=1}^{m}\mathcal{V}_i$ by Proposition \ref{prop:internal-direct-sum-family-criterion}(ii),(i). For $i\in\{1,\ldots,m\}$ and $\mathbf{x}\in\mathcal{V}_i$, $T\mathbf{x}=\lambda_i\mathbf{x}$ by idempotence and pairwise annihilation.\\
Thus $\mathcal{V}_i\subseteq E_{\lambda_i}(T)$. $\mathcal{V}_i\neq\{\mathbf{0}\}$ and $\lambda_i\in\operatorname{spec}(T)$ by $\rho_i\neq0$ and $\mathcal{V}_i\subseteq E_{\lambda_i}(T)$.\\
Conversely, if $T\mathbf{x}=\mu\mathbf{x}$ with $\mathbf{x}\neq\mathbf{0}$, write $\mathbf{x}=\sum\limits_{i=1}^{m}\mathbf{x}_i$ with $\mathbf{x}_i\in\mathcal{V}_i$.\\
Then $(\forall i\in\{1,\ldots,m\})\,(\lambda_i-\mu)\mathbf{x}_i=\mathbf{0}$ by directness of the decomposition. Since some $\mathbf{x}_i\neq\mathbf{0}$, $\mu\in\Lambda$.\\
Thus $\operatorname{spec}(T)=\Lambda$. The eigenspaces span $\mathcal{V}$ by $\mathcal{V}_i\subseteq E_{\lambda_i}(T)$ and the direct-sum decomposition.\\
Hence $T$ is diagonalizable by Theorem \ref{thm:conditions-equivalent-to-diagonalizability}(iv),(i).
\item[$(ii)$] Fix a family satisfying the right side of part (i), set $\mathcal{V}_j\defeq\operatorname{range}(\rho_j)$ for $j\in\{1,\ldots,m\}$, and fix $i\in\{1,\ldots,m\}$.\\
By the reverse direction of part (i), $\mathcal{V}=\bigoplus\limits_{j=1}^{m}\mathcal{V}_j$ and $\mathcal{V}_i\subseteq E_{\lambda_i}(T)$. If $\mathbf{x}\in E_{\lambda_i}(T)$, write $\mathbf{x}=\sum\limits_{j=1}^{m}\mathbf{x}_j$ with $\mathbf{x}_j\in\mathcal{V}_j$.\\
Then $(\forall j\in\{1,\ldots,m\}\setminus\{i\})\,\mathbf{x}_j=\mathbf{0}$ by directness of $\sum\limits_{j=1}^{m}(\lambda_j-\lambda_i)\mathbf{x}_j=\mathbf{0}$ and distinctness.\\
Hence $\mathbf{x}=\mathbf{x}_i\in\mathcal{V}_i$, and $\mathcal{V}_i=E_{\lambda_i}(T)$ by the two inclusions.
\item[$(iii)$] Fix a family satisfying the right side of part (i), set $\mathcal{V}_j\defeq\operatorname{range}(\rho_j)$, and fix $i\in\{1,\ldots,m\}$.\\
For $\mathbf{x}\in\mathcal{V}$, $\mathbf{x}=\sum\limits_{j=1}^{m}\mathbf{x}_j$ with $\mathbf{x}_j\in\mathcal{V}_j$ by reverse direction of part (i). Then $\rho_i\mathbf{x}=\mathbf{x}_i$ by idempotence, pairwise annihilation.\\
Hence $\mathbf{x}\in\ker(\rho_i)\Longleftrightarrow\mathbf{x}_i=\mathbf{0}$, so $\ker(\rho_i)=\bigoplus\limits_{j\in\{1,\ldots,m\}\setminus\{i\}}\mathcal{V}_j$. The stated identity holds by part (ii).
\item[$(iv)$] Fix a family satisfying the right side of part (i). For $i\in\{1,\ldots,m\}$, set $\ell_i(z)\defeq\prod\limits_{j\in\{1,\ldots,m\}\setminus\{i\}}\frac{z-\lambda_j}{\lambda_i-\lambda_j}$.\\
If $k\in\{1,\ldots,m\}$ and $\mathbf{x}\in E_{\lambda_k}(T)$, then $\ell_i(T)\mathbf{x}=\ell_i(\lambda_k)\mathbf{x}=\delta_{ik}\mathbf{x}=\rho_i\mathbf{x}$ by part (ii), idempotence, and pairwise annihilation.\\
The eigenspaces span $\mathcal{V}$ by parts (i)--(ii), so $\ell_i(T)=\rho_i$ by equality on a spanning family.
\end{enumerate}
\hlt{Structural Commentary}: Direct-sum coordinates construct the idempotents; Lagrange interpolation recovers them from $T$.\\
\hlt{Key Insight}: Lagrange polynomials isolate the individual eigenspaces.
\end{proof}
\end{theorem}

\begin{lemma} \label{lem:diagonalizable-restriction-to-invariant-subspace} \hlt{Diagonalizable Restrictions to Invariant Subspaces}\\
Let $\mathcal{V}$ be a nonzero finite-dimensional vector space, let $T\in\LL(\mathcal{V})$ be diagonalizable, and let $\mathcal{U}\in\operatorname{Inv}(T)$.\\
For $\lambda\in\operatorname{spec}(T)$, set $\mathcal{U}_{\lambda}\defeq\mathcal{U}\cap E_{\lambda}(T)$. The following hold:
\begin{enumerate}[label=(\roman*)]
\setlength{\itemsep}{0pt}
\item $\mathcal{U}=\bigoplus\limits_{\lambda\in\operatorname{spec}(T)}\mathcal{U}_{\lambda}$.
\item The restriction $T|_{\mathcal{U}}$ is diagonalizable.
\end{enumerate}
\begin{proof}
The family of polynomial spectral projections $(\rho_{\lambda})_{\lambda\in\operatorname{spec}(T)}$ exists by Theorem \ref{thm:algebraic-spectral-resolution-diagonalizable}(i),(iv); fix it.\\
The inclusion $T^0(\mathcal{U})=\mathcal{U}$ is immediate. If $T^k(\mathcal{U})\subseteq\mathcal{U}$, then $T^{k+1}(\mathcal{U})\subseteq T(\mathcal{U})\subseteq\mathcal{U}$ by $T$-invariance.\\
Thus $(\forall k\in\mathbb{Z}_{\geq0})\,T^k(\mathcal{U})\subseteq\mathcal{U}$ by induction, and $\rho_{\lambda}(\mathcal{U})\subseteq\mathcal{U}$ by its polynomial formula.
\begin{enumerate}[label=(\roman*)]
\setlength{\itemsep}{0pt}
\item If $\mathbf{u}\in\mathcal{U}$, then $\mathbf{u}=\sum\limits_{\lambda\in\operatorname{spec}(T)}\rho_{\lambda}\mathbf{u}$ by Theorem \ref{thm:algebraic-spectral-resolution-diagonalizable}(i).\\
The membership $\rho_{\lambda}\mathbf{u}\in\mathcal{U}_{\lambda}$ holds by the preceding inclusion and Theorem \ref{thm:algebraic-spectral-resolution-diagonalizable}(ii).\\
Hence $\mathcal{U}=\sum\limits_{\lambda\in\operatorname{spec}(T)}\mathcal{U}_{\lambda}$. The sum is direct by Proposition \ref{prop:eigenspaces-for-distinct-eigenvalues-form-direct-sum}(i).
\item If $\mathcal{U}=\{\mathbf{0}\}$, the empty basis diagonalizes $T|_{\mathcal{U}}$. Suppose $\mathcal{U}\neq\{\mathbf{0}\}$.\\
Concatenating bases of the nonzero summands forms a basis of $\mathcal{U}$ by part (i) and Theorem \ref{thm:direct-sum-iff-dimensions-add}(ii).\\
Its members are eigenvectors of $T|_{\mathcal{U}}$, so the restriction is diagonalizable by Theorem \ref{thm:conditions-equivalent-to-diagonalizability}(ii),(i).
\end{enumerate}
\hlt{Structural Commentary}: Polynomial spectral projections preserve $\mathcal{U}$ and split it across the ambient eigenspaces.\\
\hlt{Key Insight}: $T$-invariance propagates to every polynomial in $T$.
\end{proof}
\end{lemma}

\begin{theorem} \label{thm:simultaneous-diagonalization-commuting-diagonalizable-family} \hlt{Simultaneous Diagonalization of Commuting Diagonalizable Operators}\\
Let $\mathcal{V}$ be a nonzero finite-dimensional vector space, and let $\emptyset\neq\mathcal{F}\subseteq\LL(\mathcal{V})$ be finite with each member diagonalizable.\\
Then $(\forall S,T\in\mathcal{F})\,ST=TS$ iff $(\exists\text{ a basis }\beta\text{ of }\mathcal{V})$ such that $(\forall T\in\mathcal{F})\,[T]_{\beta\leftarrow\beta}$ is diagonal.
\begin{proof}
\begin{enumerate}[label=\roman*.]
\setlength{\itemsep}{0pt}
\item[$(\Rightarrow)$] Assume $(\forall S,T\in\mathcal{F})\,ST=TS$, and proceed by strong induction on $n\defeq\dim(\mathcal{V})$.\\
If $n=1$, every member of $\mathcal{F}$ is scalar, so any basis is common. Suppose $n>1$ and result holds in every smaller positive dimension.\\
If all members of $\mathcal{F}$ are scalar, any basis is common. Otherwise choose a nonscalar $S\in\mathcal{F}$.\\
By Theorem \ref{thm:conditions-equivalent-to-diagonalizability}(i),(iv), $\mathcal{V}=\bigoplus\limits_{\lambda\in\operatorname{spec}(S)}E_{\lambda}(S)$.\\
If this decomposition had one summand, then $\mathcal{V}=E_{\lambda}(S)$ and $S=\lambda I_{\mathcal{V}}$, contrary to the choice of $S$.\\
Hence $(\forall\lambda\in\operatorname{spec}(S))\,1\leq\dim E_{\lambda}(S)<n$. Fix $\lambda\in\operatorname{spec}(S)$.\\
If $T\in\mathcal{F}$ and $\mathbf{x}\in E_{\lambda}(S)$, then $S(T\mathbf{x})=T(S\mathbf{x})=\lambda T\mathbf{x}$ by commutation. Thus $E_{\lambda}(S)\in\operatorname{Inv}(T)$.\\
Set $\mathcal{F}_{\lambda}\defeq\{T|_{E_{\lambda}(S)}:T\in\mathcal{F}\}$. This family is finite, nonempty, and commuting.\\
Its members are diagonalizable by Lemma \ref{lem:diagonalizable-restriction-to-invariant-subspace}(ii). By induction, $E_{\lambda}(S)$ has a common eigenbasis for $\mathcal{F}_{\lambda}$.\\
Concatenating these bases over $\lambda\in\operatorname{spec}(S)$ forms a basis of $\mathcal{V}$ by Theorem \ref{thm:direct-sum-iff-dimensions-add}(ii).\\
The concatenated basis diagonalizes every member of $\mathcal{F}$.
\item[$(\Leftarrow)$] Choose a basis $\beta$ satisfying the right side. For $S,T\in\mathcal{F}$, set $D_S\defeq[S]_{\beta\leftarrow\beta}$ and $D_T\defeq[T]_{\beta\leftarrow\beta}$.\\
As $D_S$ and $D_T$ are diagonal, $[ST]_{\beta\leftarrow\beta}=D_SD_T=D_TD_S=[TS]_{\beta\leftarrow\beta}$ by Proposition \ref{prop:composition-and-matrix-multiplication}; $ST=TS$ by Theorem \ref{thm:linear-maps-and-matrices-same-data}(i).
\end{enumerate}
\hlt{Structural Commentary}: A nonscalar member supplies proper common invariant eigenspaces for dimensional induction.\\
\hlt{Key Insight}: Commutation preserves every eigenspace in the diagonalizable splitting.
\end{proof}
\end{theorem}

\begin{definition} \label{def:matrix-commutant} \hlt{Commutant}\\
Let $n\geq1$ and $\emptyset\neq\mathcal{F}\subseteq\operatorname{Mat}(n,n,\mathbb{F})$. The \hlt{commutant} of $\mathcal{F}$ is $\operatorname{Comm}(\mathcal{F})\defeq\{X\in\operatorname{Mat}(n,n,\mathbb{F}):(\forall A\in\mathcal{F})\,XA=AX\}$.\\
For a singleton family $\{A\}$, write $\operatorname{Comm}(A)$.
\end{definition}

\begin{lemma} \label{lem:block-diagonal-commutant-distinct-scalars} \hlt{Commuting with a Block Scalar Diagonal Matrix}\\
Let $k,n_1,\ldots,n_k\geq1$, $n\defeq\sum\limits_{i=1}^{k}n_i$, and $\lambda_1,\ldots,\lambda_k\in\mathbb{F}$. Set $\Lambda\defeq\lambda_1 I_{n_1}\oplus\cdots\oplus\lambda_k I_{n_k}$.\\
Let $A=[A_{ij}]\in\operatorname{Mat}(n,n,\mathbb{F})$ be partitioned conformally, where $A_{ij}\in\operatorname{Mat}(n_i,n_j,\mathbb{F})$.
\begin{enumerate}[label=(\roman*)]
\setlength{\itemsep}{0pt}
\item If $A$ is block diagonal, then $A\Lambda=\Lambda A$.
\item If $A\Lambda=\Lambda A$ and $\lambda_1,\ldots,\lambda_k$ are distinct, then $A$ is block diagonal.
\end{enumerate}
\begin{proof}
By block multiplication, $(A\Lambda)_{ij}=\lambda_jA_{ij}$ and $(\Lambda A)_{ij}=\lambda_iA_{ij}$.
\begin{enumerate}[label=(\roman*)]
\setlength{\itemsep}{0pt}
\item If $A$ is block diagonal, then $(\forall i,j\in\{1,\ldots,k\})\,[i\neq j\Rightarrow A_{ij}=0]$. The two diagonal-block expressions coincide, so $A\Lambda=\Lambda A$.
\item If $A\Lambda=\Lambda A$, then $(\forall i,j\in\{1,\ldots,k\})\,(\lambda_j-\lambda_i)A_{ij}=0$.\\
The relation $(\forall i,j\in\{1,\ldots,k\})\,[i\neq j\Rightarrow A_{ij}=0]$ holds by distinctness, so $A$ is block diagonal.
\end{enumerate}
\hlt{Structural Commentary}: Block comparison reduces commutation to scalar off-diagonal equations.\\
\hlt{Key Insight}: Distinct block scalars force every off-diagonal block to vanish.
\end{proof}
\end{lemma}

\begin{counterexample} \label{cex:repeated-block-scalars-do-not-force-block-diagonality} \hlt{Repeated Block Scalars Do Not Force Block Diagonality}\\
Set $k=2$, $n_1=n_2=1$, $\lambda_1=\lambda_2=1$, and $A=\left[\begin{smallmatrix}0&1\\0&0\end{smallmatrix}\right]$. Then $\Lambda=I_2$, so $A\Lambda=\Lambda A$, but $A$ is not block diagonal.\\
Thus the distinctness hypothesis in Lemma \ref{lem:block-diagonal-commutant-distinct-scalars}(ii) is necessary.
\end{counterexample}

\begin{theorem} \label{thm:commutant-of-diagonalizable-matrix} \hlt{Commutant of a Diagonalizable Matrix}\\
Let $n\geq1$, $A\in\operatorname{Mat}(n,n,\mathbb{C})$ be diagonalizable. List its distinct eigenvalues as $\mu_1,\ldots,\mu_d$, where $d\geq1$, set $n_i\defeq\dim_{\mathbb{C}}E_{\mu_i}(A)$.\\
Choose $S\in\operatorname{Mat}(n,n,\mathbb{C})^{\times}$ satisfying $A=S(\mu_1I_{n_1}\oplus\cdots\oplus\mu_dI_{n_d})S^{-1}$. Set $\mathcal{D}\defeq\{X_1\oplus\cdots\oplus X_d:(X_i)_{i=1}^{d}\in\prod\limits_{i=1}^{d}\operatorname{Mat}(n_i,n_i,\mathbb{C})\}$. Then:
\begin{enumerate}[label=(\roman*)]
\setlength{\itemsep}{0pt}
\item $(\forall B\in\operatorname{Mat}(n,n,\mathbb{C}))\,[B\in\operatorname{Comm}(A)\Longleftrightarrow S^{-1}BS\in\mathcal{D}]$.
\item $\dim_{\mathbb{C}}\operatorname{Comm}(A)=\sum\limits_{i=1}^{d}n_i^2$.
\end{enumerate}
\begin{proof}
Set $\Lambda\defeq\mu_1I_{n_1}\oplus\cdots\oplus\mu_dI_{n_d}$.
\begin{enumerate}[label=(\roman*)]
\setlength{\itemsep}{0pt}
\item[$(i)$] ($\Rightarrow$) Suppose $B\in\operatorname{Comm}(A)$. Then $(S^{-1}BS)\Lambda=\Lambda(S^{-1}BS)$ by multiplying $BA=AB$ on the left by $S^{-1}$ and on the right by $S$.\\
Thus $S^{-1}BS\in\mathcal{D}$ by Lemma \ref{lem:block-diagonal-commutant-distinct-scalars}(ii).
\item[$(i)$] ($\Leftarrow$) Suppose $S^{-1}BS\in\mathcal{D}$. Then $S^{-1}BS$ commutes with $\Lambda$ by Lemma \ref{lem:block-diagonal-commutant-distinct-scalars}(i).\\
Thus $BA=AB$ by conjugating this equality by $S$, so $B\in\operatorname{Comm}(A)$.
\item[$(ii)$] Conjugation by $S$ is a linear bijection from $\mathcal{D}$ onto $\operatorname{Comm}(A)$ by part (i), with inverse conjugation by $S^{-1}$.\\
Embedding the canonical matrix-unit bases of the $d$ diagonal blocks forms a basis of $\mathcal{D}$.\\
Hence $\dim_{\mathbb{C}}\operatorname{Comm}(A)=\dim_{\mathbb{C}}\mathcal{D}=\sum\limits_{i=1}^{d}n_i^2$ by Proposition \ref{prop:matrix-spaces-are-vector-spaces}(ii).
\end{enumerate}
\hlt{Structural Commentary}: Conjugation identifies the commutant with freely chosen endomorphisms on the eigenspaces.\\
\hlt{Key Insight}: Commuting matrices preserve eigenspaces but act freely within each one.
\end{proof}
\end{theorem}

\begin{theorem} \label{thm:commuting-family-common-eigenvector} \hlt{Common Eigenvector for a Commuting Family}\\
Let $\mathcal{V}$ be a nonzero finite-dimensional complex vector space, and let $\emptyset\neq\mathcal{F}\subseteq\LL(\mathcal{V})$ satisfy $(\forall S,T\in\mathcal{F})\,ST=TS$.\\
Then $(\exists\,(\mathbf{v},(\lambda_T)_{T\in\mathcal{F}})\in(\mathcal{V}\setminus\{\mathbf{0}\})\times\mathbb{C}^{\mathcal{F}})\,(\forall T\in\mathcal{F})\,T\mathbf{v}=\lambda_T\mathbf{v}$.
\begin{proof}
\begin{enumerate}[label=(\roman*)]
\setlength{\itemsep}{0pt}
\item \hlt{Finite family}: Suppose $\mathcal{F}=\{T_1,\ldots,T_m\}$, and induct on $m\geq1$. The case $m=1$ holds by Theorem \ref{thm:complex-operator-has-eigenvalue}. Suppose $m>1$.\\
Choose $\lambda\in\operatorname{spec}(T_1)$ by Theorem \ref{thm:complex-operator-has-eigenvalue}, and set $\mathcal{U}\defeq E_{\lambda}(T_1)\neq\{\mathbf{0}\}$.\\
If $\mathbf{u}\in\mathcal{U}$ and $j\in\{2,\ldots,m\}$, then $(T_1-\lambda I_{\mathcal{V}})T_j\mathbf{u}=T_j(T_1-\lambda I_{\mathcal{V}})\mathbf{u}=\mathbf{0}$ by commutation.\\
Thus $\mathcal{U}$ is invariant under $T_2,\ldots,T_m$. The restrictions to $\mathcal{U}$ commute and form a nonempty family of cardinality at most $m-1$.\\
They have a common eigenvector $\mathbf{v}\in\mathcal{U}\setminus\{\mathbf{0}\}$ by induction, and $\mathbf{v}$ is an eigenvector of $T_1$ by $\mathbf{v}\in E_{\lambda}(T_1)$.
\item \hlt{Finite reduction}: Set $\mathcal{W}\defeq\operatorname{span}(\mathcal{F})$. This space is finite-dimensional by Theorem \ref{thm:linear-maps-and-matrices-same-data}(ii).\\
If $\mathcal{W}=\{0\}$, then $(\forall T\in\mathcal{F})\,T=0$ by $\mathcal{F}\subseteq\mathcal{W}$. Thus any $\mathbf{v}\in\mathcal{V}\setminus\{\mathbf{0}\}$ is common with eigenvalues $\lambda_T=0$. Suppose $\mathcal{W}\neq\{0\}$.\\
Choose a basis $(R_1,\ldots,R_q)$ of $\mathcal{W}$ by Corollary \ref{cor:finite-dimensional-spaces-and-subspaces-have-bases}. Each $R_j$ is a finite linear combination of members of $\mathcal{F}$.\\
The union of their finite supports is a finite nonempty set $\mathcal{F}_0\subseteq\mathcal{F}$ satisfying $\operatorname{span}(\mathcal{F}_0)=\mathcal{W}$.
\item \hlt{Spanning extension}: Suppose $\mathcal{W}\neq\{0\}$ and let $\mathcal{F}_0$ be as in part (ii).\\
By part (i), choose $\mathbf{v}\in\mathcal{V}\setminus\{\mathbf{0}\}$ and $(\mu_T)_{T\in\mathcal{F}_0}$ satisfying $(\forall T\in\mathcal{F}_0)\,T\mathbf{v}=\mu_T\mathbf{v}$.\\
For $S\in\mathcal{F}\subseteq\mathcal{W}$, write $S=\sum\limits_{T\in\mathcal{F}_0}a_TT$ and define $\lambda_S\defeq\sum\limits_{T\in\mathcal{F}_0}a_T\mu_T$.\\
Then $S\mathbf{v}=\lambda_S\mathbf{v}$, and $\lambda_S$ is independent of the representation by $\mathbf{v}\neq\mathbf{0}$.
\end{enumerate}
\hlt{Structural Commentary}: Finite induction finds a common eigenline; a finite spanning subfamily handles arbitrary $\mathcal{F}$.\\
\hlt{Key Insight}: Commutation preserves the first operator's chosen eigenspace.
\end{proof}
\end{theorem}

\begin{theorem} \label{thm:simultaneous-schur-triangularization-commuting-family} \hlt{Simultaneous Schur Triangularization}\\
Let $n\geq1$ and let $\emptyset\neq\mathcal{F}\subseteq\operatorname{Mat}(n,n,\mathbb{C})$ satisfy $(\forall A,B\in\mathcal{F})\,AB=BA$. The following hold:
\begin{enumerate}[label=(\roman*)]
\setlength{\itemsep}{0pt}
\item $(\exists\text{ a unitary }U\in\operatorname{Mat}(n,n,\mathbb{C}))$ such that $(\forall A\in\mathcal{F})\,U^*AU$ is upper triangular.
\item Suppose $\mathcal{F}\subseteq\operatorname{Mat}(n,n,\mathbb{R})$ and $(\forall A\in\mathcal{F})\,\operatorname{spec}(A)\subseteq\mathbb{R}$.\\
Then $(\exists\text{ an orthogonal }U\in\operatorname{Mat}(n,n,\mathbb{R}))$ such that $(\forall A\in\mathcal{F})\,U^tAU$ is upper triangular.
\end{enumerate}
\begin{proof}
Proceed by strong induction on $n$. Both conclusions hold with $U=I_1$ when $n=1$.\\
Suppose $n>1$ and both conclusions hold in each smaller positive dimension.
\begin{enumerate}[label=(\roman*)]
\setlength{\itemsep}{0pt}
\item By Theorem \ref{thm:commuting-family-common-eigenvector}, the operators on $\mathbb{C}^n$ have a common eigenvector.\\
Normalize it to $\mathbf{x}$ and extend it to a unitary matrix $V=[\mathbf{x}\ V_2]$ by Theorem \ref{thm:extend-isometric-columns-to-unitary-matrix}(ii).\\
For $A\in\mathcal{F}$, write $M_A\defeq V^*AV=\left[\begin{smallmatrix}\lambda_A&\mathbf{r}_A\\0&\widehat A\end{smallmatrix}\right]$. The lower-left block vanishes by $A\mathbf{x}=\lambda_A\mathbf{x}$.\\
If $A,B\in\mathcal{F}$, then $M_AM_B=V^*ABV=V^*BAV=M_BM_A$, so $\widehat A\widehat B=\widehat B\widehat A$ by comparison of the lower-right blocks.\\
Thus $\widehat{\mathcal{F}}\defeq\{\widehat A:A\in\mathcal{F}\}$ is nonempty and commuting.\\
By induction, a unitary $W$ simultaneously upper triangularizes its members. Set $U\defeq V(I_1\oplus W)$.\\
The matrix $U$ is unitary by Proposition \ref{prop:algebra-of-unitary-matrices}(i)--(ii), and $(\forall A\in\mathcal{F})\,U^*AU$ is upper triangular.
\item Regard $\mathcal{F}$ as a commuting family on $\mathbb{C}^n$. A common eigenvector $\mathbf{z}\in\mathbb{C}^n\setminus\{\mathbf{0}\}$ exists by Theorem \ref{thm:commuting-family-common-eigenvector}; fix it.\\
Fix $(\lambda_A)_{A\in\mathcal{F}}$ satisfying $(\forall A\in\mathcal{F})\,A\mathbf{z}=\lambda_A\mathbf{z}$, and write $\mathbf{z}=\mathbf{u}+i\mathbf{v}$.\\
The relation $(\forall A\in\mathcal{F})\,\lambda_A\in\operatorname{spec}(A)\subseteq\mathbb{R}$ holds by the eigenvector identities and spectral hypothesis.\\
Because each $A$ is real, $A\mathbf{u}=\lambda_A\mathbf{u}$ and $A\mathbf{v}=\lambda_A\mathbf{v}$ by comparison of real and imaginary parts.\\
At least one of $\mathbf{u},\mathbf{v}$ is nonzero, so a common real unit eigenvector $\mathbf{x}$ exists by normalization.\\
Extend $\mathbf{x}$ to a real orthogonal matrix $V=[\mathbf{x}\ V_2]$ by Theorem \ref{thm:extend-isometric-columns-to-unitary-matrix}(ii).\\
For $A\in\mathcal{F}$, write $M_A\defeq V^tAV=\left[\begin{smallmatrix}\lambda_A&\mathbf{r}_A\\0&\widehat A\end{smallmatrix}\right]$. The lower-right blocks are real and commute by the calculation in part (i).\\
Fix $A\in\mathcal{F}$ and $\mu\in\operatorname{spec}(\widehat A)$, and choose $\mathbf{y}\in\mathbb{C}^{n-1}\setminus\{\mathbf{0}\}$ satisfying $\widehat A\mathbf{y}=\mu\mathbf{y}$.\\
If $\mu=\lambda_A$, then $\mu\in\mathbb{R}$. If $\mu\neq\lambda_A$, set $c\defeq\frac{\mathbf{r}_A\mathbf{y}}{\mu-\lambda_A}$ and $\mathbf{w}\defeq(c,\mathbf{y})^t$.\\
Then $\mathbf{w}\neq\mathbf{0}$ and $M_A\mathbf{w}=(\mu c,\mu\mathbf{y})^t=\mu\mathbf{w}$. Since $V\mathbf{w}\neq\mathbf{0}$ by orthogonality of $V$, it is an eigenvector of $A$ for $\mu$.\\
Hence $\mu\in\mathbb{R}$ by the spectral hypothesis. Thus $(\forall A\in\mathcal{F})\,\operatorname{spec}(\widehat A)\subseteq\mathbb{R}$.\\
By induction, a real orthogonal $W$ simultaneously upper triangularizes $\widehat{\mathcal{F}}$.\\
Then $U\defeq V(I_1\oplus W)$ is real orthogonal by Proposition \ref{prop:algebra-of-unitary-matrices}(i)--(ii), and $(\forall A\in\mathcal{F})\,U^tAU$ is upper triangular.
\end{enumerate}
\hlt{Structural Commentary}: A common eigenline creates the first block, and commutation descends to the compression.\\
\hlt{Key Insight}: Real spectra turn a complex common eigenvector into a real one.
\end{proof}
\end{theorem}

Figure \ref{fig:simultaneous-triangularization-commuting} depicts the complex induction in Theorem \ref{thm:simultaneous-schur-triangularization-commuting-family}.\\
A common eigenline creates the first block, and pairwise commutation passes to the lower-right blocks.\\
The theorem, rather than the diagram, owns the simultaneous triangularization conclusion.

\begin{illustrationfigure}[H]
  \centering
  \begin{tikzpicture}[
  >=Stealth,
  every node/.style={font=\small},
  box/.style={draw=gray!55, fill=gray!5, rounded corners=2pt, align=center, inner xsep=5pt, inner ysep=4pt},
  hi/.style={draw=blue!60!black, fill=blue!7, rounded corners=2pt, align=center, inner xsep=5pt, inner ysep=4pt},
  arr/.style={->, line width=0.75pt, blue!70!black}
]
  \node[hi] (common) at (0,0)
    {$\mathbb{C}\mathbf{x}$};
  \node[box] (block) at (3.18,0)
    {$V^*AV=\begin{bmatrix}\lambda_A&*\\0&A'\end{bmatrix}$};
  \node[box] (comm) at (6.60,0)
    {$A'B'=B'A'$};

  \draw[arr] (common.east) -- (block.west);
  \draw[arr] (block.east) -- (comm.west);
\end{tikzpicture}

  \caption{A shared eigenline determines a simultaneous first block; lower-right compressions remain commuting.}
  \label{fig:simultaneous-triangularization-commuting}
\end{illustrationfigure}

\begin{corollary} \label{cor:commuting-sum-product-eigenvalues} \hlt{Eigenvalues of Sums and Products of Commuting Matrices}\\
Let $n\geq1$ and let $A,B\in\operatorname{Mat}(n,n,\mathbb{C})$ satisfy $AB=BA$. Then $(\exists\,((\lambda_i)_{i=1}^{n},(\mu_i)_{i=1}^{n})\in(\mathbb{C}^n)^2)$ such that the following hold:
\begin{enumerate}[label=(\roman*)]
\setlength{\itemsep}{0pt}
\item $\operatorname{spec}(A)=\{\lambda_1,\ldots,\lambda_n\}$ and $\operatorname{spec}(B)=\{\mu_1,\ldots,\mu_n\}$;
\item $\operatorname{spec}(A+B)=\{\lambda_i+\mu_i:i\in\{1,\ldots,n\}\}$;
\item $\operatorname{spec}(AB)=\{\lambda_i\mu_i:i\in\{1,\ldots,n\}\}$.
\end{enumerate}
\begin{proof}
Choose a unitary $U$ by Theorem \ref{thm:simultaneous-schur-triangularization-commuting-family}(i). Set $T\defeq U^*AU$, $R\defeq U^*BU$, $\lambda_i\defeq t_{ii}$, $\mu_i\defeq r_{ii}$. Fix $M\in\operatorname{Mat}(n,n,\mathbb{C})$ and $\theta\in\mathbb{C}$.\\
If $\mathbf{x}\neq\mathbf{0}$, then $(U^*MU)\mathbf{x}=\theta\mathbf{x}\Longleftrightarrow M(U\mathbf{x})=\theta(U\mathbf{x})$ and $U\mathbf{x}\neq\mathbf{0}$ by unitarity of $U$. Hence $\operatorname{spec}(U^*MU)=\operatorname{spec}(M)$.
\begin{enumerate}[label=(\roman*)]
\setlength{\itemsep}{0pt}
\item The matrices $T$ and $R$ are upper triangular with diagonals $(\lambda_i)_{i=1}^{n}$ and $(\mu_i)_{i=1}^{n}$.\\
Their spectra are underlying sets of these lists by Proposition \ref{prop:eigenvalues-of-upper-triangular-operator}. Identities in part (i) hold by preceding similarity calculation.
\item The matrix $T+R=U^*(A+B)U$ is upper triangular with diagonal $(\lambda_i+\mu_i)_{i=1}^{n}$.\\
Its spectrum is the asserted set by Proposition \ref{prop:eigenvalues-of-upper-triangular-operator}. The identity in part (ii) holds by the preceding similarity calculation.
\item The matrix $TR=U^*ABU$ is upper triangular with diagonal $(\lambda_i\mu_i)_{i=1}^{n}$.\\
Its spectrum is the asserted set by Proposition \ref{prop:eigenvalues-of-upper-triangular-operator}. The identity in part (iii) holds by the preceding similarity calculation.
\end{enumerate}
\hlt{Structural Commentary}: Common triangular coordinates reduce both spectra to diagonal arithmetic.\\
\hlt{Key Insight}: A common diagonal pairing controls both resulting spectra.
\end{proof}
\end{corollary}
